\documentclass[10pt,twocolumn]{article}

\usepackage{amsmath,amssymb}
\usepackage{graphicx}
\usepackage{booktabs}
\usepackage{xcolor}
\usepackage{colortbl}
\usepackage{geometry}
\usepackage{microtype}
\usepackage{caption}
\usepackage{hyperref}
\usepackage{titlesec}

\geometry{top=2cm, bottom=2cm, left=1.5cm, right=1.5cm, columnsep=0.6cm}

\definecolor{bestrow}{RGB}{198,239,206}   % green highlight
\definecolor{seccolor}{RGB}{30,80,160}

\titleformat{\section}{\normalfont\bfseries\color{seccolor}}{{\thesection.}}{0.5em}{}
\titleformat{\subsection}{\normalfont\itshape}{}{0em}{}

\captionsetup{font=small, labelfont=bf}

% -----------------------------------------------------------------------
\title{\textbf{Speed Control of a BDC Motor with Vertical Rod:\\
From First-Principles Dynamics to Data-Driven Feedforward}}
\author{Paweekorn Buasakorn}
\date{}
% -----------------------------------------------------------------------

\begin{document}
\maketitle
\thispagestyle{empty}

% ===================================================================
\section{Motor Dynamics}
% ===================================================================

A brushed DC (BDC) motor driving a vertical rigid rod is modelled
through two coupled subsystems.

\subsection{Electrical subsystem}
Applying Kirchhoff's voltage law to the armature circuit:
\begin{equation}
  L\,\frac{di}{dt} = V_{\mathrm{in}} - Ri - K_b\,\omega
  \label{eq:elec}
\end{equation}
where $V_{\mathrm{in}}$ is the applied voltage, $i$ the armature
current, $\omega$ the shaft angular speed, $R$ the armature
resistance, $L$ the inductance, and $K_b$ the back-EMF constant.

\subsection{Mechanical subsystem}
The net torque balance on the motor shaft with an attached rod of
mass $m$, length $\ell$, centre of mass at $\ell_{cm}=\ell/2$, and
moment of inertia $I_{\mathrm{rod}}=\tfrac{1}{3}m\ell^{2}$ gives:
\begin{equation}
  J_{\mathrm{eff}}\,\frac{d\omega}{dt}
    = K_t\,i - B\,\omega - mg\ell_{cm}\sin\theta
  \label{eq:mech}
\end{equation}
where $J_{\mathrm{eff}}=J_m + I_{\mathrm{rod}}$, $K_t$ the torque
constant, $B$ the viscous friction coefficient, $\theta$ the rod angle
from the downward vertical, and $g$ gravitational acceleration.
The system is integrated numerically using a fixed-step RK4 integrator
with $\Delta t = 10^{-3}$\,s.

The parameters used throughout are listed in Table~\ref{tab:params}.

\begin{table}[h]
\centering
\caption{Simulation parameters}
\label{tab:params}
\begin{tabular}{llll}
\toprule
\textbf{Symbol} & \textbf{Value} & \textbf{Symbol} & \textbf{Value}\\
\midrule
$R$   & 3.0\,$\Omega$       & $J_m$  & $7.4\times10^{-5}$\,kg\,m$^2$\\
$L$   & 4\,mH               & $B$    & $5\times10^{-3}$\,N\,m\,s/rad\\
$K_t$ & 0.05\,N\,m/A        & $m$    & 0.05\,kg\\
$K_b$ & 0.05\,V\,s/rad      & $\ell$ & 0.1\,m\\
$V_{\max}$ & 12\,V          & $g$    & 9.81\,m/s$^2$\\
\bottomrule
\end{tabular}
\end{table}

% ===================================================================
\section{Gravity Compensation at Steady State}
% ===================================================================

To derive a feedforward term for gravity, consider the steady-state
condition where all derivatives vanish
($\frac{di}{dt}=0$,\; $\frac{d\omega}{dt}=0$):
\begin{align}
  V &= Ri_{ss} + K_b\,\omega \label{eq:ss_elec}\\
  0 &= K_t\,i_{ss} - B\,\omega - mg\ell_{cm}\sin\theta \label{eq:ss_mech}
\end{align}
Solving \eqref{eq:ss_mech} for $i_{ss}$ and substituting into
\eqref{eq:ss_elec}:
\begin{equation}
  V_{ss} = \underbrace{\frac{R\,mg\ell_{cm}\sin\theta}{K_t}}_{V_{\mathrm{grav}}}
           + \underbrace{\left(\frac{RB}{K_t}+K_b\right)\omega}_{V_{\omega}}
  \label{eq:vss}
\end{equation}
The gravity feedforward term $V_{\mathrm{grav}}$ is extracted and added
directly to the PI control output:
\begin{equation}
  V_{\mathrm{cmd}} = V_{\mathrm{grav}} + K_p\,e + K_i\!\int\! e\,dt
  \label{eq:pi_grav}
\end{equation}

% ===================================================================
\section{Inverse Model via Laplace Transform}
% ===================================================================

Taking the Laplace transform of \eqref{eq:elec} and \eqref{eq:mech}
and eliminating $I(s)$:
\begin{equation}
  V_{\mathrm{in}}(s)
    = \underbrace{\frac{(J_{\mathrm{eff}}s+B)(Ls+R)+K_tK_b}{K_t}}_{G_1(s)}\,\Omega(s)
    + \underbrace{\frac{Ls+R}{K_t}}_{G_2(s)}\,T_{\mathrm{load}}(s)
  \label{eq:laplace_inv}
\end{equation}
where $T_{\mathrm{load}}(s) = mg\ell_{cm}\sin\theta$ is treated as a
disturbance. The term $G_1(s)$ gives the voltage required to produce a
desired speed trajectory $\Omega_{\mathrm{ref}}(s)$, while $G_2(s)$
accounts for load rejection.

% ===================================================================
\section{Realizability and the Low-Pass Correction}
% ===================================================================

Expanding $G_1(s)$:
\begin{equation}
  G_1(s) = \frac{J_{\mathrm{eff}}L}{K_t}\,s^2
          + \frac{J_{\mathrm{eff}}R + BL}{K_t}\,s
          + \frac{BR + K_tK_b}{K_t}
  \label{eq:g1}
\end{equation}
This is a \emph{polynomial} in $s$ (improper transfer function): its
numerator degree (2) exceeds its denominator degree (0).
Physically, applying $G_1(s)$ requires twice-differentiating the
reference $\omega_{\mathrm{ref}}$, which amplifies noise and is
non-causal in practice.

To make the feedforward realizable, a second-order low-pass filter
$(\tau s+1)^{-2}$ is appended:
\begin{equation}
  G_{ff}(s) = \frac{G_1(s)}{(\tau s+1)^2}
  \label{eq:gff}
\end{equation}
The filter introduces two poles and shifts relative degree to zero,
making $G_{ff}(s)$ proper. However, choosing $\tau$ requires careful
tuning and the resulting filter still demands precise knowledge of all
motor parameters ($R,L,K_t,K_b,J,B$), making it fragile and
impractical outside of MATLAB/Simulink environments.

% ===================================================================
\section{Polynomial Regression as a Practical Alternative}
% ===================================================================

\subsection{Motivation}
Rather than implementing the model-based inverse filter, we adopt a
\emph{data-driven} approach: measure the steady-state relationship
between voltage and speed experimentally, invert it, and fit a
polynomial.

\subsection{Data collection}
The motor is driven at each voltage level
$V \in \{-12, -9, \ldots, 9, 12\}$\,V.
The simulation waits for steady state (rolling 1\,s window with
$\sigma_\omega < 0.5$\,rad/s and drift $<0.3$\,rad/s), then records
$N=100$ noisy speed samples per level, yielding a dataset
$\{(V_k,\,\omega_k)\}$.

\subsection{Inversion and zero-bias fit}
Because the motor is symmetric and exhibits no offset at $\omega=0$,
the inverse map $V = f(\omega)$ is constrained to pass through the
origin. A zero-intercept polynomial of degree $n$ is fitted by solving:
\begin{equation}
  \min_{\mathbf{c}}\,\bigl\|X\mathbf{c} - \bar{V}\bigr\|^2,\quad
  X_{ij} = \bar{\omega}_i^{\,j},\;\; j=1,\ldots,n
  \label{eq:lstsq}
\end{equation}
using \texttt{numpy.linalg.lstsq}, where $\bar{\omega}_i$ and
$\bar{V}_i$ are the per-level means. With $n=2$ the result is:
\begin{equation}
  \hat{V}(\omega) = c_1\,\omega + c_2\,\omega^2
  \label{eq:poly}
\end{equation}
The fitted coefficients are $c_1 \approx 0.3588$\,V\,s/rad and
$c_2 \approx 0$, consistent with the theoretical steady-state gain
$\tfrac{RB}{K_t}+K_b = 0.35$\,V\,s/rad. At runtime the feedforward
voltage is simply:
\begin{equation}
  V_{ff}^{\mathrm{poly}}(\omega_{\mathrm{ref}})
    = \hat{V}\!\left(\operatorname{clip}(\omega_{\mathrm{ref}},\,
      \omega_{\min},\,\omega_{\max})\right)
  \label{eq:poly_ff}
\end{equation}
This requires \emph{no parameter knowledge}, is trivially fast to
evaluate, and generalises to any actuator with a monotone
steady-state characteristic.

% ===================================================================
\section{Experimental Comparison}
% ===================================================================

Eight closed-loop experiments compare controllers on two reference
profiles: a constant step ($\omega_{\mathrm{ref}}=10$\,rad/s) and a
quintic trajectory that smoothly ramps from 0 to 10\,rad/s between
$t=0.5$\,s and $t=3.5$\,s with zero boundary velocity and
acceleration. All controllers share $K_p=5$, $K_i=2$.

Table~\ref{tab:results} reports the RMS tracking error over 5\,s.
The best result is highlighted in green.

\begin{table}[h]
\centering
\caption{RMS tracking error (rad/s) for each controller and reference.}
\label{tab:results}
\begin{tabular}{lcc}
\toprule
\textbf{Controller} & \textbf{Const.\ ref} & \textbf{Trajectory}\\
\midrule
PI only                        & 0.495 & 0.322 \\
PI + gravity FFW               & 0.450 & 0.279 \\
PI + poly FFW                  & 0.376 & 0.171 \\
\rowcolor{bestrow}
PI + poly FFW + gravity FFW    & 0.322 & \textbf{0.033} \\
\bottomrule
\end{tabular}
\end{table}

\subsection{Discussion}
Pure PI achieves basic regulation but carries a steady-state offset
under gravity load. Adding the analytical gravity feedforward
(Section~2) reduces the error marginally on constant
references but still leaves the integral to compensate model mismatch.

The polynomial feedforward provides a larger benefit: by pre-computing
the steady-state voltage for any $\omega_{\mathrm{ref}}$, it removes
the dominant tracking lag and allows the PI to act only on residual
disturbances.

Combining both feedforwards yields the best performance. On the quintic
trajectory the RMSE drops to \textbf{0.033\,rad/s} --- an order of
magnitude below pure PI --- because the polynomial term handles the
voltage demand and the gravity term compensates the angle-dependent
disturbance, leaving the integrator nearly idle. The constant-reference
experiments show smaller absolute differences because the integrator
eventually winds up to eliminate any steady bias, narrowing the gap
between strategies.

% ===================================================================
\section{Conclusion}
% ===================================================================

A complete simulation pipeline was developed for a BDC motor driving a
vertical rod.  The model-based inverse feedforward, while theoretically
exact, is non-realizable without parameter knowledge and additional
filtering.  Polynomial regression on steady-state input--output data
provides a practical, model-free substitute that, combined with an
analytical gravity term, reduces trajectory-tracking error by
${\sim}10\times$ compared to PI alone. The pipeline --- data
collection, curve fitting, controller design, and benchmarking --- is
fully automated and reproducible from the project root.

\end{document}
